%=============================================================================
% Wave 4, Iteration 18: QUASI-STATIC REGIME PROOF
% Agent 12 (Cosmology --- Quasi-Static Regime Proof)
%
% Model: Opus 4.6
% Date: 20 March 2026
%
% INHERITED STATE (from i17):
%   - c_s^2 = (W'/K'') * c^2/4 ~ c^2 at ALL epochs (i17 Finding 2)
%   - Psi-field CANNOT cluster on sub-Hubble scales (i17 confirmed)
%   - i17 identified 3 scenarios: (a) c_s ~ c FATAL, (b) c_s -> 0 viable,
%     (c) quasi-static -> psi slaved to matter, "c_s" irrelevant
%   - G_eff(k) = G_N(k^2 + k_mu^2)/k^2 with k_mu ~ 0.01 Mpc^{-1}
%   - i17 found P(k) deficit ~10^{-6} on small scales (k >> k_mu)
%   - i17 concluded "catastrophic deficit" --- escape routes 1--6 all fail
%   - Skordis & Zlosnik (2021) AeST theory reproduces CMB + P(k) via
%     ghost condensate + quasi-static scalar
%
% MANDATE (i18):
%   Prove or disprove scenario (c): quasi-static regime.
%   Is the psi-field slaved to matter such that c_s is irrelevant?
%   Compare with AQUAL/RMOND/AeST literature.
%   TOP 5 findings.
%=============================================================================

\documentclass[11pt]{article}
\usepackage{amsmath,amssymb,amsthm}
\usepackage{geometry}
\geometry{margin=1in}
\usepackage{booktabs}
\usepackage{enumitem}
\usepackage{hyperref}
\usepackage{xcolor}
\usepackage{tcolorbox}
\usepackage{float}
\usepackage{siunitx}

\newtheorem{theorem}{Theorem}
\newtheorem{proposition}{Proposition}
\newtheorem{finding}{Finding}
\newtheorem{correction}{Correction}
\newtheorem{result}{Result}
\newtheorem{lemma}{Lemma}

\newcommand{\ee}[1]{\times 10^{#1}}
\newcommand{\psifield}{\psi}
\newcommand{\DFD}{\textsc{dfd}}
\newcommand{\GR}{\textsc{gr}}
\newcommand{\LCDM}{$\Lambda$CDM}
\newcommand{\Geff}{G_{\rm eff}}
\newcommand{\aM}{a_0}
\newcommand{\Hzero}{H_0}
\newcommand{\kmu}{k_\mu}

\title{\textbf{Wave 4, Iteration 18:}\\
\textbf{Quasi-Static Regime Proof}\\[12pt]
{\large Is the DFD $\psi$-Field Slaved to Matter? Does $c_s$ Become Irrelevant?}\\[4pt]
{\large Comparison with AQUAL, RMOND/AeST, and k-Essence Literature}\\[4pt]
{\large Five Findings with Full Derivation Chains}}
\author{DFD Research Programme --- W4i18 Agent 12 (Opus 4.6)}
\date{20 March 2026}

\begin{document}
\maketitle
\tableofcontents
\newpage

%=============================================================================
\section*{Executive Summary: Top 5 Findings}
%=============================================================================

\begin{tcolorbox}[colback=blue!5!white, colframe=blue!70!black,
  title=\textbf{W4i18 TOP 5 FINDINGS --- QUASI-STATIC REGIME PROOF}]

\begin{enumerate}

\item \textbf{FINDING 1 (QUASI-STATIC REGIME IS VALID FOR ALL SUB-HUBBLE
  MODES --- BUT THIS DOES NOT SAVE DFD):
  The temporal term $K'' \cdot (c/a_0)^2 \ddot{\delta\psi}$ IS negligible
  compared to the spatial term $W' \cdot |\nabla\delta\psi|^2/a_*^2$
  for ALL modes with $k \gg aH/c$. The ratio
  temporal/spatial $\sim H^2/k^2 c_s^2 \ll 1$ for sub-Hubble modes
  regardless of the value of $c_s$.
  The psi-field satisfies an effective Poisson equation
  $\nabla^2 \delta\psi = \text{source}$ at each instant.
  HOWEVER, this quasi-static slaving does NOT provide CDM-like
  potential wells because the quasi-static response IS the $\Geff$
  enhancement already computed in i17, and $\Geff \approx G_N$
  on small scales ($k \gg k_\mu$).}

  \textbf{Derivation}: The linearized perturbation equation from
  Eq.~(2.9) of section02\_formalism.tex is:
  \begin{equation}
  K''(\Delta_{\rm bg}) \frac{c^2}{a_0^2} \ddot{\delta\psi}
  - \frac{W'(y_{\rm bg})}{a_*^2} \nabla^2 \delta\psi
  = \frac{4\pi G}{c^2} \delta\rho.
  \end{equation}
  For a mode with physical wavenumber $k$ and frequency $\omega$:
  \begin{equation}
  \frac{\text{temporal}}{\text{spatial}}
  = \frac{K'' (c/a_0)^2 \omega^2}{W'/a_*^2 \cdot k^2}
  = \frac{\omega^2}{c_s^2 k^2},
  \end{equation}
  where $c_s^2 = (W'/K'') \cdot c^2/4$ (from i17).

  For gravitationally-driven modes: $\omega \sim H$ (modes grow on
  Hubble timescale, not oscillate). For sub-Hubble modes: $k \gg H/c$.
  Therefore:
  \begin{equation}
  \frac{\text{temporal}}{\text{spatial}}
  \sim \frac{H^2}{c_s^2 (H/c)^2}
  = \frac{c^2}{c_s^2} \cdot \frac{H^2}{k^2 c^2}
  \sim \frac{c^2}{c_s^2} \cdot \left(\frac{aH}{k c}\right)^2.
  \end{equation}

  Since $c_s \sim c/2$ (from i17): $c^2/c_s^2 \sim 4$.
  For sub-Hubble modes ($k \gg aH/c$): the factor $(aH/kc)^2 \ll 1$.
  Therefore: temporal/spatial $\sim 4 \times (aH/kc)^2 \ll 1$.

  \textbf{The quasi-static approximation is valid for ALL sub-Hubble
  modes, at ALL epochs.} The large $c_s$ actually HELPS the
  quasi-static limit hold: the faster the propagation speed, the
  more rapidly the field adjusts to its quasi-static configuration,
  and the more negligible the temporal term becomes relative to
  the spatial gradient.

  \textbf{CRITICAL INSIGHT}: Having $c_s \sim c$ does NOT mean the
  field fails to respond to matter. It means the field responds
  INSTANTANEOUSLY (on sub-Hubble timescales) to the matter
  distribution. The field is ``slaved'' to matter precisely because
  $c_s$ is large, not despite it. A field with $c_s \to 0$ would be
  slow to respond and would NOT be quasi-static.

  \textbf{The quasi-static equation is}:
  \begin{equation}
  \nabla \cdot \left[\mu\!\left(\frac{|\nabla\psi|}{a_*}\right)
  \nabla\psi\right] = -\frac{8\pi G}{c^2} \delta\rho.
  \label{eq:QS-poisson}
  \end{equation}
  This is the SPATIAL sector of the DFD field equation, with the
  temporal sector dropped. It is a nonlinear Poisson equation.

\item \textbf{FINDING 2 (THE i17 ``CATASTROPHIC DEFICIT'' IS CORRECT
  AND QUASI-STATIC DOES NOT RESOLVE IT):
  The quasi-static psi response to baryon perturbations IS the
  $\Geff(k)$ enhancement. Linearizing Eq.~(\ref{eq:QS-poisson})
  gives $\Geff(k) = G_N (k^2 + k_\mu^2)/k^2$, which provides
  negligible enhancement on small scales. The baryon growth
  equation $\ddot\delta_b + 2H\dot\delta_b = 4\pi \Geff(k)
  \bar\rho_b \delta_b$ with $\Geff \approx G_N$ for $k \gg k_\mu$
  yields $P_{\rm DFD}(k)/P_{\rm CDM}(k) \sim 10^{-6}$
  at $k > 0.1\;\text{Mpc}^{-1}$. The i17 result stands.}

  The quasi-static scenario (c) and the $c_s \sim c$ scenario (a)
  are NOT alternatives --- they are THE SAME THING viewed from
  different angles:
  \begin{itemize}[nosep]
  \item Scenario (a): $c_s \sim c$ means psi-perturbations propagate
    at speed $c$ and are smoothed on sub-Hubble scales. Therefore
    the psi-field does not independently cluster.
  \item Scenario (c): The psi-field is quasi-static (slaved to matter)
    because it adjusts on timescale $\sim 1/(c_s k) \ll 1/H$.
    The field tracks the matter distribution instantaneously but
    does not add independent clustering power.
  \end{itemize}

  These are equivalent statements. The quasi-static psi-field provides
  a modified gravitational response (captured by $\Geff$) but does
  NOT act as an independently clustering component. The ``dark matter''
  role requires an independently clustering component with its own
  perturbation amplitude built up from $z_{\rm eq} \sim 3400$.

  \textbf{Quantitative confirmation}: The baryon growth from
  $z_{\rm dec} = 1100$ to $z = 0$ under $\Geff = G_N$ gives
  $\delta_b(0) \sim 10^{-5} \times 1100 \sim 10^{-2}$. In \LCDM,
  CDM has $\delta_{\rm CDM}(z_{\rm dec}) \sim 10^{-2}$ (from growth
  during $z_{\rm eq}$ to $z_{\rm dec}$), and baryons catch up to
  $\delta_{\rm CDM} \sim 10$ by $z = 0$. DFD baryons never reach
  the nonlinear regime on galaxy scales.

\item \textbf{FINDING 3 (COMPARISON WITH AQUAL/AeST --- THE SKORDIS-ZLOSNIK
  SOLUTION AND WHY DFD LACKS IT):
  Skordis \& Zlosnik (2021) achieved a viable CMB + $P(k)$ in their
  AeST (Aether-Scalar-Tensor) theory NOT through a quasi-static scalar
  alone, but through a ``ghost condensate'' vector field whose
  perturbations provide an independently clustering dark-matter-like
  component. The scalar is quasi-static; the VECTOR provides
  the CDM-like wells. DFD has only the scalar. This is the
  structural reason DFD cannot reproduce $P(k)$.}

  The AeST theory (Skordis \& Zlosnik 2021, PRL 127, 161302) contains:
  \begin{enumerate}[nosep]
  \item A scalar field $\varphi$ (MOND-like, quasi-static in galaxies)
  \item A unit timelike vector field $A^\mu$ (``aether'')
  \item A free function $\mathcal{K}$ controlling the vector kinetics
  \end{enumerate}

  The crucial ingredient is the vector field $A^\mu$. On cosmological
  scales, its perturbations $\delta A^\mu$ behave as a pressureless
  fluid ($c_s = 0$) that clusters from $z_{\rm eq}$ onward, exactly
  mimicking CDM. The scalar $\varphi$ is quasi-static and provides
  the MOND limit in galaxies. The vector provides the CDM limit in
  cosmology. This two-field structure is essential.

  \textbf{DFD has only one scalar field $\psi$.} There is no vector
  (or tensor) degree of freedom that can provide CDM-like clustering.
  The scalar is quasi-static (as proven in Finding 1), which means it
  tracks matter but does not independently cluster. There is no second
  field to fill the CDM role.

  \textbf{Comparison with original AQUAL} (Bekenstein \& Milgrom 1984):
  AQUAL was a non-relativistic theory and did not address cosmological
  perturbation growth. Its relativistic extensions (TeVeS, Bekenstein 2004)
  introduced additional fields --- a vector and a scalar --- precisely
  to handle the cosmological clustering problem. The Skordis-Zlosnik
  AeST is the latest and most successful in this lineage. All viable
  relativistic MOND theories that match $P(k)$ require at least TWO
  field types: a scalar for MOND galaxies and a vector/tensor for
  CDM-like cosmological clustering.

  \textbf{Implication for DFD}: A single scalar field cannot
  simultaneously be quasi-static (giving MOND) and independently
  clustering (giving CDM). These are contradictory requirements.
  DFD would need a second degree of freedom --- either modifying
  the vector sector, promoting the TT sector, or introducing a
  new field --- to achieve viable cosmology.

\item \textbf{FINDING 4 (THE k-ESSENCE QUASI-STATIC LITERATURE CONFIRMS
  THE PROBLEM --- OSCILLATING MODES ARE SUPPRESSED):
  In k-essence cosmology, the quasi-static approximation yields
  a matter density perturbation almost identical to \LCDM{} only
  when the k-essence field serves as dark energy (background),
  not as dark matter (perturbations). When k-essence perturbations
  are expected to cluster like matter, the oscillating solution
  (neglected in the quasi-static limit) becomes important, and
  its suppression by $c_s \sim c$ prevents clustering.
  DFD's psi-field is structurally analogous to k-essence with
  $c_s \sim c/2$, which places it firmly in the ``dark energy''
  category, not ``dark matter.''}

  From Bamba \& Matsumoto (2012, PRD 85, 084026): the quasi-static
  solution for matter perturbations in k-essence is nearly identical
  to \LCDM{} if the background is tuned. However, there exists an
  \textit{oscillating solution} that is neglected in the quasi-static
  approximation. For k-essence with $c_s \sim c$, these oscillations
  have wavelength $\sim c/H$ (Hubble scale) and are damped on
  sub-Hubble scales. The k-essence field cannot cluster below its
  sound horizon $\lambda_s = c_s/H$.

  For DFD: $c_s \sim c/2$, so $\lambda_s \sim c/(2H) \sim d_H/2$
  (half the Hubble radius). The psi-field is smooth on ALL sub-Hubble
  scales. This is identical to quintessence-type dark energy with
  $c_s = c$, which is known to contribute to the expansion rate
  but not to structure formation.

  \textbf{The Sawicki-Bellini criterion} (2015): The quasi-static
  approximation is valid inside the dark field's sound horizon. For
  DFD, the sound horizon is $\sim d_H/2$, so the QSA is valid on
  all sub-Hubble scales (confirming Finding 1). But validity of the
  QSA means precisely that the field does not independently cluster
  --- it is slaved to the matter distribution.

  \textbf{For comparison}: CDM has $c_s = 0$, so its sound horizon
  is zero. The QSA is never valid for CDM perturbations (they evolve
  dynamically, not quasi-statically). This dynamic evolution is what
  allows CDM to grow potential wells independently of baryons.

\item \textbf{FINDING 5 (THE i15 ``$c_s^2 \to 0$ THEOREM'' IS MISLEADING
  BUT NOT WRONG --- IT DESCRIBES THE BACKGROUND EoS, NOT PERTURBATION
  PROPAGATION. SCENARIO (c) IS CONFIRMED BUT DOES NOT HELP.
  THE THREE SCENARIOS COLLAPSE TO ONE):}

  Resolving the three scenarios from i17:
  \begin{itemize}[nosep]
  \item \textbf{Scenario (a)} [$c_s \sim c$, FATAL]: CONFIRMED.
    The perturbation propagation speed is $c_s = c/2$ from the kinetic
    matrix. Psi-perturbations do not cluster on sub-Hubble scales.

  \item \textbf{Scenario (b)} [$c_s \to 0$, viable]: WRONG as stated.
    The Appendix Q ``$c_s^2 \to 0$'' result refers to the adiabatic
    sound speed $dp/d\rho$ of the \textit{background} psi-fluid,
    not to the perturbation propagation speed. The background has
    $w \to 0$ (dust-like EoS) in the $\Delta \to 0$ limit, but the
    perturbation propagation speed remains $c_s = c/2$. These are
    distinct quantities in k-essence theories.

  \item \textbf{Scenario (c)} [quasi-static, $c_s$ irrelevant]:
    CONFIRMED but DOES NOT HELP. The quasi-static limit IS valid
    (Finding 1). The psi-field IS slaved to matter. But ``slaved
    to matter'' means it provides only $\Geff$ enhancement, which
    is the same as scenarios (a) and (c). The three scenarios
    are not alternatives --- they are the same physics described
    differently.
  \end{itemize}

  \textbf{THE THREE SCENARIOS COLLAPSE TO ONE}: The psi-field has
  $c_s \sim c/2$, is quasi-static on sub-Hubble scales, responds
  instantaneously to matter, enhances gravity via $\Geff(k)$, but
  does NOT independently cluster. All three descriptions are
  equivalent. The matter power spectrum deficit is real, robust, and
  not resolvable within the single-scalar DFD framework.

  \textbf{STRUCTURAL DIAGNOSIS}: DFD's cosmological viability requires
  one of:
  \begin{enumerate}[nosep]
  \item A second clustering degree of freedom (vector field, as in AeST)
  \item A mechanism to make $K''$ large enough that $c_s \ll c$ for
    perturbations (requires modifying the temporal kinetic function
    beyond the simple $\mu$-interpolation)
  \item A demonstration that baryonic physics alone (Silk damping,
    reionization, nonlinear feedback) can compensate the $10^{-6}$
    deficit (implausible given the magnitude)
  \item Acceptance that DFD is a galactic-scale theory only, not a
    cosmological theory (as original MOND/AQUAL were)
  \end{enumerate}

  Option 1 (AeST-like extension) is the most promising path but
  requires extending DFD beyond its current single-scalar structure.
  Option 2 requires $K''(\Delta_{\rm bg}) \gg 1$, i.e., $K$ must be
  extremely convex at the cosmological background value. For
  $c_s^2 \sim c^2/4 \cdot W'/K''$ to give $c_s \ll c$, we need
  $K'' \gg 1$ at $\Delta_{\rm bg} \sim 1.5$. The current simple
  $\mu$ function gives $K''(1.5) = 2/(1+1.5)^3 = 0.128$, yielding
  $c_s^2 = c^2/4 \times 1/0.128 = 2c^2$ (superluminal). A viable
  DFD cosmology would require $K''(\Delta) \gtrsim 10^6$ to bring
  $c_s$ below the BAO scale, which is grossly incompatible with
  the smooth $\mu(x) = x/(1+x)$ interpolation.

\end{enumerate}
\end{tcolorbox}


%=============================================================================
%=============================================================================
\part{Finding 1: Quasi-Static Validity Proof}
\label{part:QS-proof}
%=============================================================================
%=============================================================================

\section{The DFD Perturbation Equation}

From section02\_formalism.tex, the full dynamic DFD action (Eq.~2.9) is:
\begin{equation}
S_\psi = \int dt\,d^3x\;\left\{
  \frac{a_*^2}{8\pi G}\left[
    W\!\left(\frac{|\nabla\psi|^2}{a_*^2}\right)
    + K\!\left(\frac{c}{a_0}|\dot\psi{-}\dot\psi_0|\right)
  \right]
  - \frac{c^2}{2}\,\psi(\rho - \bar\rho)
\right\}.
\end{equation}

Linearizing around the cosmological background $\bar\psi$,
$\psi = \bar\psi + \delta\psi$, the perturbation equation is:
\begin{equation}
\underbrace{K''(\Delta_{\rm bg}) \frac{c^2}{a_0^2}}_{\text{temporal coefficient}}
\ddot{\delta\psi}
- \underbrace{\frac{W'(y_{\rm bg})}{a_*^2}}_{\text{spatial coefficient}}
\nabla^2 \delta\psi
= \frac{4\pi G}{c^2} \delta\rho.
\label{eq:pert-full}
\end{equation}

\subsection{The Temporal-to-Spatial Ratio}

For a Fourier mode with comoving wavenumber $k$ and frequency $\omega$:
\begin{equation}
\mathcal{R}(k,\omega) \equiv
\frac{K''(c/a_0)^2 \omega^2}{W'/a_*^2 \cdot k^2}
= \frac{\omega^2}{c_s^2 k^2},
\end{equation}
where $c_s^2 = (W'/K'') \cdot a_0^2/(a_*^2 c^2)
= (W'/K'') \cdot c^2/4$ (using $a_* = 2a_0/c^2$).

\subsection{Gravitationally-Driven Modes}

For modes driven by gravitational instability (not freely oscillating):
$\omega \sim H$ (growth rate set by Hubble expansion).
For sub-Hubble modes: $k \gg aH/c$.

Therefore:
\begin{equation}
\mathcal{R} \sim \frac{H^2}{c_s^2 k^2}
= \frac{H^2}{(c^2/4) \cdot (W'/K'') \cdot k^2}.
\end{equation}

At the BAO scale $k_{\rm BAO} \sim 0.06\;\text{Mpc}^{-1}$ and $z = 0$:
$H_0/(c k_{\rm BAO}) \sim 2.2\ee{-18}/(3\ee{8} \times 0.06/3.086\ee{22})
= 2.2\ee{-18}/5.8\ee{-16} = 3.8\ee{-3}$.
So $\mathcal{R} \sim 4 \times (3.8\ee{-3})^2 = 5.8\ee{-5} \ll 1$.

At galaxy scales $k \sim 1\;\text{Mpc}^{-1}$:
$\mathcal{R} \sim 4 \times (2.3\ee{-4})^2 = 2\ee{-7} \ll 1$.

\begin{theorem}[Quasi-Static Validity]
\label{thm:QS}
The temporal term in the DFD perturbation equation is negligible
compared to the spatial term for all modes satisfying $k \gg aH/c$
(sub-Hubble). The ratio is:
\begin{equation}
\mathcal{R}(k) = \frac{4}{W'/K''} \cdot \frac{H^2}{c^2 k^2} \ll 1.
\end{equation}
Since $W'/K'' \geq 1$ for all DFD interpolating functions and all
epochs, the quasi-static approximation holds for all sub-Hubble modes
at all cosmological times.
\end{theorem}

\subsection{Why Large $c_s$ Helps, Not Hurts}

A common misconception: $c_s \sim c$ should mean the field ``propagates
away'' and cannot respond to matter. The opposite is true.

The quasi-static approximation works BETTER when $c_s$ is large:
\begin{itemize}[nosep]
\item Large $c_s$: the field adjusts to changes in the source on
  timescale $\tau_{\rm adj} \sim 1/(c_s k) \ll 1/H$. The field
  tracks the source nearly instantaneously.
\item Small $c_s$: the field adjusts slowly, $\tau_{\rm adj} \sim
  1/(c_s k) \sim 1/H$ or longer. The temporal term is NOT negligible.
  The QSA breaks down.
\item $c_s = 0$: the field does not propagate at all. It evolves
  dynamically (temporal and spatial terms comparable). This is the
  CDM regime --- dynamical clustering, not quasi-static slaving.
\end{itemize}

The DFD psi-field with $c_s \sim c/2$ is in the maximally quasi-static
regime. It is perfectly slaved to the baryon distribution at every
instant.


%=============================================================================
%=============================================================================
\part{Finding 2: Quasi-Static Does Not Resolve the $P(k)$ Deficit}
\label{part:Pk-deficit}
%=============================================================================
%=============================================================================

\section{The Quasi-Static Baryon Growth Equation}

Dropping the temporal term from Eq.~(\ref{eq:pert-full}):
\begin{equation}
\frac{W'(y_{\rm bg})}{a_*^2} \nabla^2 \delta\psi
= -\frac{4\pi G}{c^2} \delta\rho_b.
\end{equation}

In the Newtonian regime ($y_{\rm bg} \gg 1$, $W' \approx 1$):
$\nabla^2 \delta\psi = -4\pi G a_*^2 \delta\rho_b/c^2
= -(8\pi G/c^2) \cdot (a_0/c^2)^2 \delta\rho_b$.

Wait --- let me be more careful. The full quasi-static equation is
the nonlinear Poisson (Eq.~\ref{eq:QS-poisson}):
$\nabla \cdot [\mu(|\nabla\psi|/a_*) \nabla\psi]
= -(8\pi G/c^2)(\rho_b - \bar\rho_b)$.

Linearizing $\psi = \bar\psi + \delta\psi$, and noting that the
background $\bar\psi$ is homogeneous (no spatial gradients on the
cosmological background):
\begin{equation}
\mu'(0) \cdot \frac{|\nabla\delta\psi|}{a_*} \cdot \nabla^2\delta\psi
+ \mu(0) \cdot \nabla^2\delta\psi = -\frac{8\pi G}{c^2}\delta\rho_b.
\end{equation}

Since $\mu(0) = 0$ and $\mu'(0) = 1$ (for the simple function):
$\nabla^2 \delta\psi / a_* \cdot |\nabla\delta\psi| = -(8\pi G/c^2)\delta\rho_b$.

This is the MOND regime equation. For a Fourier mode:
$-k^2 \delta\psi \cdot k|\delta\psi|/a_* = -(8\pi G/c^2)\delta\rho_b$,
which gives $|\delta\psi|^2 = (8\pi G a_*)/(c^2 k^3) \delta\rho_b$.

This nonlinear relation makes Fourier analysis non-trivial. The
standard procedure (as in AQUAL literature) is to define an effective
$\Geff(k)$ through linearization around the LOCAL background gradient.
On cosmological scales, the relevant background gradient is
$|\nabla\bar\psi|/a_* \sim a_0/(c^2 a_*)$ from the Hubble flow,
which gives $\mu \sim 1$ (Newtonian regime) on scales where
$|\nabla\bar\psi|/a_* \gg 1$.

\section{The $\Geff(k)$ from Quasi-Static Response}

From the i17 derivation (which used the linearized quasi-static
equation around the MOND background):
\begin{equation}
\Geff(k) = G_N \cdot \frac{k^2 + \kmu^2}{k^2},
\quad \kmu \approx 0.01\;\text{Mpc}^{-1}.
\end{equation}

The baryon growth equation:
\begin{equation}
\ddot\delta_b + 2H\dot\delta_b = 4\pi \Geff(k) \bar\rho_b \delta_b.
\end{equation}

At $k = 0.1\;\text{Mpc}^{-1}$ (cluster scale): $\Geff/G_N = 1.01$.
At $k = 1\;\text{Mpc}^{-1}$ (galaxy scale): $\Geff/G_N = 1.0001$.

The growth factor from $z = 1100$ to $z = 0$ in the matter era:
$D(z) \propto a \propto (1+z)^{-1}$, so
$\delta_b(0) \sim \delta_b(1100) \times 1100 \sim 10^{-5} \times 1100
= 1.1\ee{-2}$.

In \LCDM: $\delta_{\rm CDM+b}(0) \sim 10$ (nonlinear on these scales).

\textbf{Deficit: $P_{\rm DFD}/P_{\rm CDM} \sim (10^{-2}/10)^2 = 10^{-6}$
on scales $k > 0.1\;\text{Mpc}^{-1}$.}

This confirms the i17 result: the quasi-static interpretation does
not change the conclusion. The $\Geff$ enhancement from the quasi-static
psi-field is precisely what was already computed, and it is insufficient.


%=============================================================================
%=============================================================================
\part{Finding 3: The AeST/RMOND Lesson}
\label{part:AeST}
%=============================================================================
%=============================================================================

\section{Why AeST Succeeds Where DFD Fails}

The Skordis-Zlosnik AeST theory (2021) achieves:
\begin{enumerate}[nosep]
\item CMB power spectrum matching \LCDM{} to $<1\%$
\item Matter power spectrum matching \LCDM{} including BAO
\item MOND limit in quasi-static galactic environments
\item GW speed $c_T = c$ (GW170817 compatible)
\item No ghost instabilities
\end{enumerate}

The structural ingredients:
\begin{enumerate}[nosep]
\item \textbf{Metric tensor $g_{\mu\nu}$} (standard GR graviton)
\item \textbf{Scalar field $\varphi$} (MOND-like, quasi-static in galaxies)
\item \textbf{Unit timelike vector $A^\mu$} (``aether'')
\item \textbf{Free function $\mathcal{K}$} controlling vector kinetics
\end{enumerate}

The vector field $A^\mu$ is the key. In the cosmological background,
$A^\mu = (1, 0, 0, 0)$ in comoving coordinates (``ghost condensate'').
Perturbations $\delta A^i$ behave as an irrotational pressureless
fluid --- exactly like CDM. The ``ghost condensate density'' evolves as
$\rho_{\rm gc} \propto a^{-3}$, clusters from $z_{\rm eq}$ onward,
and provides the CDM potential wells into which baryons fall after
decoupling.

The scalar $\varphi$ is quasi-static in the galactic environment,
providing the MOND modification. It does NOT cluster cosmologically.

\section{Structural Comparison: AeST vs DFD}

\begin{center}
\begin{tabular}{lcc}
\toprule
Feature & AeST & DFD \\
\midrule
MOND limit (galaxies) & Scalar $\varphi$ (QS) & Scalar $\psi$ (QS) \\
CDM-like clustering & Vector $A^\mu$ & \textbf{NONE} \\
GW speed $= c$ & Yes & Yes \\
Background $H(z)$ & \LCDM-like & \LCDM{} (by construction) \\
CMB $C_\ell$ & Matches & Matches (smooth $\psi$) \\
$P(k)$ & Matches & \textbf{$10^{-6}$ deficit} \\
DOF count & Scalar + Vector + Tensor & Scalar + Tensor \\
\bottomrule
\end{tabular}
\end{center}

\textbf{The lesson is clear}: A quasi-static scalar can provide MOND
in galaxies. But cosmological structure formation requires an additional
clustering degree of freedom. In TeVeS, this was a vector. In AeST,
it is a refined vector (ghost condensate). In DFD, no such field exists.

\section{Could DFD's TT Sector Help?}

DFD has a TT gravitational wave sector with $c_T = c$. Could tensor
perturbations provide clustering?

No: tensor perturbations (gravitational waves) propagate at $c$ and
decay as $1/a$. They do not cluster. The TT sector is radiative, not
matter-like.

\section{Could a Vector Extension Save DFD?}

A DFD extension with an additional unit timelike vector $A^\mu$ could,
in principle, mimic the AeST structure. The vector perturbations would
cluster like CDM on large scales while the scalar $\psi$ provides MOND
in galaxies. However:
\begin{itemize}[nosep]
\item This would change DFD from a single-scalar theory to a
  scalar-vector theory
\item The optical metric framework would need extension
\item The ``derivation from first principles'' (refractive medium)
  would need justification for the vector field
\item The parameter space would expand significantly
\end{itemize}

This is a major structural modification, not a minor fix.


%=============================================================================
%=============================================================================
\part{Finding 4: k-Essence Quasi-Static Literature}
\label{part:k-essence}
%=============================================================================
%=============================================================================

\section{k-Essence Perturbation Theory}

DFD's scalar sector is a k-essence theory: the Lagrangian depends
nonlinearly on kinetic terms. The k-essence perturbation literature
(Bamba \& Matsumoto 2012; Vikman 2007) establishes:

\begin{enumerate}[nosep]
\item The perturbation equation admits two classes of solutions:
  (a) quasi-static (growing/decaying) and (b) oscillating.
\item The oscillating modes have frequency $\omega \sim c_s k$
  and are damped by Hubble friction.
\item The quasi-static solution dominates when $c_s k \gg H$
  (all sub-Hubble modes for $c_s \sim c$).
\item The quasi-static matter perturbation in k-essence mimics
  \LCDM{} IF the k-essence serves as dark energy (smooth
  background), NOT as dark matter.
\end{enumerate}

\textbf{Application to DFD}: The psi-field has $c_s \sim c/2$,
so $c_s k \gg H$ for all sub-Hubble modes. The quasi-static
solution dominates (confirming Finding 1). The oscillating
modes are smoothed out on the Hubble scale.

The matter (baryon) perturbation under the quasi-static psi-field
is:
\begin{equation}
\ddot\delta_b + 2H\dot\delta_b
= 4\pi G_N (1 + \epsilon(k)) \bar\rho_b \delta_b,
\end{equation}
where $\epsilon(k) = \Geff/G_N - 1 = k_\mu^2/k^2$.

For $k \gg k_\mu$: $\epsilon \to 0$, recovering standard baryon
growth. This is the k-essence quasi-static result: the dark energy
field's effect on matter perturbations is negligible on small scales.

\section{The Sound Horizon as the Critical Scale}

The Sawicki-Bellini (2015) analysis of Horndeski theories establishes
that the quasi-static approximation is valid inside the dark field's
sound horizon:
\begin{equation}
k > k_{\rm sh} = \frac{aH}{c_s}.
\end{equation}

For DFD: $c_s \sim c/2$, so $k_{\rm sh} \sim 2aH/c$. All
cosmologically observable modes ($k > 10^{-4}\;\text{Mpc}^{-1}$) are
well inside the sound horizon.

\textbf{For CDM}: $c_s = 0$, so $k_{\rm sh} \to \infty$. The QSA
is NEVER valid for CDM. CDM perturbations evolve dynamically,
which is precisely what allows independent clustering.

This establishes a fundamental dichotomy:
\begin{center}
\begin{tabular}{lcc}
\toprule
Property & $c_s \sim c$ (DFD $\psi$) & $c_s = 0$ (CDM) \\
\midrule
QSA validity & All sub-Hubble & Never \\
Perturbation behavior & Slaved to source & Independent dynamics \\
Clustering & Tracks matter & Grows independently \\
Role & Modified gravity & Dark matter \\
\bottomrule
\end{tabular}
\end{center}


%=============================================================================
%=============================================================================
\part{Finding 5: Resolution of the Three Scenarios}
\label{part:resolution}
%=============================================================================
%=============================================================================

\section{The Three Scenarios Are One}

i17 identified three scenarios:
\begin{description}
\item[Scenario (a)] $c_s \sim c$ $\to$ FATAL for clustering.
\item[Scenario (b)] $c_s \to 0$ $\to$ viable (CDM-like).
\item[Scenario (c)] Quasi-static $\to$ $c_s$ irrelevant.
\end{description}

\textbf{Resolution}:
\begin{itemize}
\item Scenario (a) is CORRECT: $c_s^2 = (W'/K'') \cdot c^2/4 \sim c^2$.
  This is the perturbation propagation speed from the kinetic matrix
  of the action. It is robust, following directly from the DFD action.

\item Scenario (b) is MISLEADING: The Appendix Q ``$c_s^2 \to 0$''
  theorem refers to $dp/d\rho$ for the background psi-fluid (the
  thermodynamic adiabatic sound speed), not the perturbation
  propagation speed. In k-essence theories, these differ. The
  background EoS can have $w \to 0$ (dust-like) while the
  perturbation speed remains $c_s \sim c/2$. This is not a
  contradiction --- it is a standard feature of k-essence.

  \textbf{Physical analogy}: A cosmological constant has $w = -1$
  and $c_s^2 = -c^2$ (imaginary, meaning no propagation). A stiff
  fluid has $w = 1$ and $c_s = c$. DFD's psi-field has $w \to 0$
  (dust-like background) but $c_s = c/2$ (perturbations propagate
  fast). This combination is perfectly consistent in k-essence.

\item Scenario (c) is CORRECT but UNHELPFUL: The quasi-static limit
  IS valid (proven in Part~\ref{part:QS-proof}). The psi-field IS
  slaved to matter. But this slaving is precisely the $\Geff$
  enhancement, which is insufficient on small scales.
\end{itemize}

\section{The Structural Impossibility}

The DFD psi-field faces a structural impossibility:
\begin{enumerate}
\item To provide MOND in galaxies, $\psi$ must be quasi-static
  (respond to matter distribution via nonlinear Poisson equation).
  This requires $c_s \gtrsim c$ so the field adjusts faster than
  the dynamics.

\item To provide CDM-like clustering in cosmology, $\psi$ must have
  $c_s \to 0$ so it clusters independently of baryons.

\item Requirements 1 and 2 are contradictory. A single scalar field
  cannot have $c_s \sim c$ (for galactic MOND) and $c_s \to 0$
  (for cosmological CDM) simultaneously.
\end{enumerate}

This is not specific to DFD --- it is a general theorem about
single-scalar MOND theories. The solution adopted by AeST is to
split the two roles between different fields (scalar for MOND,
vector for CDM).

\section{Implications for the DFD Programme}

\begin{tcolorbox}[colback=red!5!white, colframe=red!70!black,
  title=\textbf{COSMOLOGICAL STATUS OF DFD (Post-i18)}]

\textbf{CONFIRMED VIABLE}:
\begin{itemize}[nosep]
\item Galactic rotation curves (MOND phenomenology)
\item Gravitational lensing ($\Geff$ in quasi-static limit)
\item Solar system (Newtonian limit, PPN parameters)
\item GW propagation ($c_T = c$)
\item CMB $C_\ell$ (smooth psi-field at $z \sim 1100$)
\item Background expansion $H(z)$ (unmodified by construction)
\end{itemize}

\textbf{FATALLY DEFICIENT}:
\begin{itemize}[nosep]
\item Matter power spectrum $P(k)$ at $k > 0.01\;\text{Mpc}^{-1}$:
  suppressed by $\sim 10^{-6}$ relative to observations
\item Galaxy formation, cluster mass function, Lyman-$\alpha$ forest:
  all require CDM-like potential wells that DFD cannot provide
\item $\sigma_8$: predicted $\sim 100\times$ too low
\end{itemize}

\textbf{POSSIBLE EXTENSIONS}:
\begin{itemize}[nosep]
\item AeST-like vector extension (most promising, but major revision)
\item Modified $K(\Delta)$ with $K'' \gg 1$ at cosmological $\Delta$
  (requires abandoning simple $\mu$ interpolation; likely introduces
  instabilities or violates galactic phenomenology)
\item Accept DFD as galactic-scale theory (historical MOND position)
\end{itemize}

\textbf{mochi\_class remains EXISTENTIALLY URGENT}: The numerical
Boltzmann code would quantify the deficit precisely, determine
whether the $k_\mu$ scale provides any partial rescue at BAO
scales, and test whether nonlinear effects modify the simple
estimates above. The qualitative conclusion ($P(k)$ deficit) is
robust, but the quantitative magnitude ($10^{-6}$ vs $10^{-3}$
vs $10^{-1}$) matters for determining whether an extension is
a minor modification or a major overhaul.

\end{tcolorbox}


%=============================================================================
\section*{References and Comparison Literature}
%=============================================================================

\begin{enumerate}[nosep]
\item Skordis \& Zlosnik (2021), PRL 127, 161302: ``New Relativistic
  Theory for Modified Newtonian Dynamics'' --- AeST theory, CMB + $P(k)$
  match via ghost condensate vector.
\item Bekenstein \& Milgrom (1984), ApJ 286, 7: AQUAL --- original
  non-relativistic modified gravity action. No cosmological perturbation
  treatment.
\item Bamba \& Matsumoto (2012), PRD 85, 084026: k-essence perturbation
  theory. Quasi-static vs oscillating solutions. Sound speed dependence.
\item Sawicki \& Bellini (2015), PRD 92, 084061: Quasi-static
  approximation validity in Horndeski theories. Sound horizon criterion.
\item Pace et al.\ (2021), JCAP 2021, 017: Comparison of quasi-static
  approaches in Horndeski models. Sub-percent accuracy inside sound
  horizon.
\item Noller et al.\ (2014), JCAP 2014, 026: Relativistic scalar fields
  and quasi-static approximation in modified gravity.
\item Bekenstein (2004), PRD 70, 083509: TeVeS --- relativistic MOND
  with vector field. Precursor to AeST.
\end{enumerate}


\end{document}
